% Options for packages loaded elsewhere
% Options for packages loaded elsewhere
\PassOptionsToPackage{unicode}{hyperref}
\PassOptionsToPackage{hyphens}{url}
\PassOptionsToPackage{dvipsnames,svgnames,x11names}{xcolor}
%
\documentclass[
  11pt,
  letterpaper,
  DIV=11,
  numbers=noendperiod]{scrartcl}
\usepackage{xcolor}
\usepackage[margin=1in]{geometry}
\usepackage{amsmath,amssymb}
\setcounter{secnumdepth}{5}
\usepackage{iftex}
\ifPDFTeX
  \usepackage[T1]{fontenc}
  \usepackage[utf8]{inputenc}
  \usepackage{textcomp} % provide euro and other symbols
\else % if luatex or xetex
  \usepackage{unicode-math} % this also loads fontspec
  \defaultfontfeatures{Scale=MatchLowercase}
  \defaultfontfeatures[\rmfamily]{Ligatures=TeX,Scale=1}
\fi
\usepackage{lmodern}
\ifPDFTeX\else
  % xetex/luatex font selection
\fi
% Use upquote if available, for straight quotes in verbatim environments
\IfFileExists{upquote.sty}{\usepackage{upquote}}{}
\IfFileExists{microtype.sty}{% use microtype if available
  \usepackage[]{microtype}
  \UseMicrotypeSet[protrusion]{basicmath} % disable protrusion for tt fonts
}{}
\makeatletter
\@ifundefined{KOMAClassName}{% if non-KOMA class
  \IfFileExists{parskip.sty}{%
    \usepackage{parskip}
  }{% else
    \setlength{\parindent}{0pt}
    \setlength{\parskip}{6pt plus 2pt minus 1pt}}
}{% if KOMA class
  \KOMAoptions{parskip=half}}
\makeatother
% Make \paragraph and \subparagraph free-standing
\makeatletter
\ifx\paragraph\undefined\else
  \let\oldparagraph\paragraph
  \renewcommand{\paragraph}{
    \@ifstar
      \xxxParagraphStar
      \xxxParagraphNoStar
  }
  \newcommand{\xxxParagraphStar}[1]{\oldparagraph*{#1}\mbox{}}
  \newcommand{\xxxParagraphNoStar}[1]{\oldparagraph{#1}\mbox{}}
\fi
\ifx\subparagraph\undefined\else
  \let\oldsubparagraph\subparagraph
  \renewcommand{\subparagraph}{
    \@ifstar
      \xxxSubParagraphStar
      \xxxSubParagraphNoStar
  }
  \newcommand{\xxxSubParagraphStar}[1]{\oldsubparagraph*{#1}\mbox{}}
  \newcommand{\xxxSubParagraphNoStar}[1]{\oldsubparagraph{#1}\mbox{}}
\fi
\makeatother

\usepackage{color}
\usepackage{fancyvrb}
\newcommand{\VerbBar}{|}
\newcommand{\VERB}{\Verb[commandchars=\\\{\}]}
\DefineVerbatimEnvironment{Highlighting}{Verbatim}{commandchars=\\\{\}}
% Add ',fontsize=\small' for more characters per line
\usepackage{framed}
\definecolor{shadecolor}{RGB}{255,255,255}
\newenvironment{Shaded}{\begin{snugshade}}{\end{snugshade}}
\newcommand{\AlertTok}[1]{\textcolor[rgb]{0.75,0.01,0.01}{\textbf{\colorbox[rgb]{0.97,0.90,0.90}{#1}}}}
\newcommand{\AnnotationTok}[1]{\textcolor[rgb]{0.79,0.38,0.79}{#1}}
\newcommand{\AttributeTok}[1]{\textcolor[rgb]{0.00,0.34,0.68}{#1}}
\newcommand{\BaseNTok}[1]{\textcolor[rgb]{0.69,0.50,0.00}{#1}}
\newcommand{\BuiltInTok}[1]{\textcolor[rgb]{0.39,0.29,0.61}{\textbf{#1}}}
\newcommand{\CharTok}[1]{\textcolor[rgb]{0.57,0.30,0.62}{#1}}
\newcommand{\CommentTok}[1]{\textcolor[rgb]{0.54,0.53,0.53}{#1}}
\newcommand{\CommentVarTok}[1]{\textcolor[rgb]{0.00,0.58,1.00}{#1}}
\newcommand{\ConstantTok}[1]{\textcolor[rgb]{0.67,0.33,0.00}{#1}}
\newcommand{\ControlFlowTok}[1]{\textcolor[rgb]{0.12,0.11,0.11}{\textbf{#1}}}
\newcommand{\DataTypeTok}[1]{\textcolor[rgb]{0.00,0.34,0.68}{#1}}
\newcommand{\DecValTok}[1]{\textcolor[rgb]{0.69,0.50,0.00}{#1}}
\newcommand{\DocumentationTok}[1]{\textcolor[rgb]{0.38,0.47,0.50}{#1}}
\newcommand{\ErrorTok}[1]{\textcolor[rgb]{0.75,0.01,0.01}{\underline{#1}}}
\newcommand{\ExtensionTok}[1]{\textcolor[rgb]{0.00,0.58,1.00}{\textbf{#1}}}
\newcommand{\FloatTok}[1]{\textcolor[rgb]{0.69,0.50,0.00}{#1}}
\newcommand{\FunctionTok}[1]{\textcolor[rgb]{0.39,0.29,0.61}{#1}}
\newcommand{\ImportTok}[1]{\textcolor[rgb]{1.00,0.33,0.00}{#1}}
\newcommand{\InformationTok}[1]{\textcolor[rgb]{0.69,0.50,0.00}{#1}}
\newcommand{\KeywordTok}[1]{\textcolor[rgb]{0.12,0.11,0.11}{\textbf{#1}}}
\newcommand{\NormalTok}[1]{\textcolor[rgb]{0.12,0.11,0.11}{#1}}
\newcommand{\OperatorTok}[1]{\textcolor[rgb]{0.12,0.11,0.11}{#1}}
\newcommand{\OtherTok}[1]{\textcolor[rgb]{0.00,0.43,0.16}{#1}}
\newcommand{\PreprocessorTok}[1]{\textcolor[rgb]{0.00,0.43,0.16}{#1}}
\newcommand{\RegionMarkerTok}[1]{\textcolor[rgb]{0.00,0.34,0.68}{\colorbox[rgb]{0.88,0.91,0.97}{#1}}}
\newcommand{\SpecialCharTok}[1]{\textcolor[rgb]{0.24,0.68,0.91}{#1}}
\newcommand{\SpecialStringTok}[1]{\textcolor[rgb]{1.00,0.33,0.00}{#1}}
\newcommand{\StringTok}[1]{\textcolor[rgb]{0.75,0.01,0.01}{#1}}
\newcommand{\VariableTok}[1]{\textcolor[rgb]{0.00,0.34,0.68}{#1}}
\newcommand{\VerbatimStringTok}[1]{\textcolor[rgb]{0.75,0.01,0.01}{#1}}
\newcommand{\WarningTok}[1]{\textcolor[rgb]{0.75,0.01,0.01}{#1}}

\usepackage{longtable,booktabs,array}
\usepackage{calc} % for calculating minipage widths
% Correct order of tables after \paragraph or \subparagraph
\usepackage{etoolbox}
\makeatletter
\patchcmd\longtable{\par}{\if@noskipsec\mbox{}\fi\par}{}{}
\makeatother
% Allow footnotes in longtable head/foot
\IfFileExists{footnotehyper.sty}{\usepackage{footnotehyper}}{\usepackage{footnote}}
\makesavenoteenv{longtable}
\usepackage{graphicx}
\makeatletter
\newsavebox\pandoc@box
\newcommand*\pandocbounded[1]{% scales image to fit in text height/width
  \sbox\pandoc@box{#1}%
  \Gscale@div\@tempa{\textheight}{\dimexpr\ht\pandoc@box+\dp\pandoc@box\relax}%
  \Gscale@div\@tempb{\linewidth}{\wd\pandoc@box}%
  \ifdim\@tempb\p@<\@tempa\p@\let\@tempa\@tempb\fi% select the smaller of both
  \ifdim\@tempa\p@<\p@\scalebox{\@tempa}{\usebox\pandoc@box}%
  \else\usebox{\pandoc@box}%
  \fi%
}
% Set default figure placement to htbp
\def\fps@figure{htbp}
\makeatother





\setlength{\emergencystretch}{3em} % prevent overfull lines

\providecommand{\tightlist}{%
  \setlength{\itemsep}{0pt}\setlength{\parskip}{0pt}}



 


\KOMAoption{captions}{tableheading}
\makeatletter
\@ifpackageloaded{caption}{}{\usepackage{caption}}
\AtBeginDocument{%
\ifdefined\contentsname
  \renewcommand*\contentsname{Table of contents}
\else
  \newcommand\contentsname{Table of contents}
\fi
\ifdefined\listfigurename
  \renewcommand*\listfigurename{List of Figures}
\else
  \newcommand\listfigurename{List of Figures}
\fi
\ifdefined\listtablename
  \renewcommand*\listtablename{List of Tables}
\else
  \newcommand\listtablename{List of Tables}
\fi
\ifdefined\figurename
  \renewcommand*\figurename{Figure}
\else
  \newcommand\figurename{Figure}
\fi
\ifdefined\tablename
  \renewcommand*\tablename{Table}
\else
  \newcommand\tablename{Table}
\fi
}
\@ifpackageloaded{float}{}{\usepackage{float}}
\floatstyle{ruled}
\@ifundefined{c@chapter}{\newfloat{codelisting}{h}{lop}}{\newfloat{codelisting}{h}{lop}[chapter]}
\floatname{codelisting}{Listing}
\newcommand*\listoflistings{\listof{codelisting}{List of Listings}}
\makeatother
\makeatletter
\makeatother
\makeatletter
\@ifpackageloaded{caption}{}{\usepackage{caption}}
\@ifpackageloaded{subcaption}{}{\usepackage{subcaption}}
\makeatother
\usepackage{bookmark}
\IfFileExists{xurl.sty}{\usepackage{xurl}}{} % add URL line breaks if available
\urlstyle{same}
\hypersetup{
  pdftitle={Technical Appendix: Sparse LD and CSR},
  pdfauthor={Peter Sørensen},
  colorlinks=true,
  linkcolor={blue},
  filecolor={Maroon},
  citecolor={Blue},
  urlcolor={Blue},
  pdfcreator={LaTeX via pandoc}}


\title{Technical Appendix: Sparse LD and CSR}
\author{Peter Sørensen}
\date{}
\begin{document}
\maketitle

\renewcommand*\contentsname{Table of contents}
{
\hypersetup{linkcolor=}
\setcounter{tocdepth}{3}
\tableofcontents
}

This appendix defines the sparse linkage disequilibrium (LD)
representation used by the current \texttt{sblr} BED-to-CSR workflow. It
describes the statistical matrix, compressed sparse row (CSR) structure,
disk-backed construction, indexing conventions, alignment requirements,
and downstream conversion used by \texttt{stblr\_csr()} and related
wrappers.

The stable user-facing contract is deliberately narrower than the
current binary-file implementation:

\begin{enumerate}
\def\labelenumi{\arabic{enumi}.}
\tightlist
\item
  create sparse LD with \texttt{make\_sparseLD()}, which records
  metadata in \texttt{Glist\$sparseLD};
\item
  inspect the resulting prefix with \texttt{sparseLD\_read\_CSR()} when
  needed; and
\item
  pass the updated \texttt{Glist} to \texttt{stblr\_csr()} or
  \texttt{stblr\_csr\_annot()}.
\end{enumerate}

\texttt{sparseLD\_stream\_CSR()} is the lower-level writer used by the
preparation helper. Its \texttt{ld\_prefix} remains the disk-backed
artifact identifier, but the current high-level workflow carries that
prefix through \texttt{Glist\$sparseLD}.

The present file suffixes, integer widths, floating-point widths,
metadata keys, and upper-triangle storage convention are documented here
to explain behavior and support validation. They should not be treated
as a guaranteed manual file-format API. Use package readers, writers,
and samplers rather than editing or independently constructing these
files.

\section{From BED Genotypes to Sufficient Statistics and Sparse
LD}\label{from-bed-genotypes-to-sufficient-statistics-and-sparse-ld}

For genotype matrix \(X\), phenotype vector \(y\), and marker effects
\(b\), a summary-statistics marker sampler needs quantities including

\[
wy=X^\top y,
\qquad
ww_i=x_i^\top x_i,
\qquad
yy=y^\top y,
\]

and enough off-diagonal marker information to update

\[
r=X^\top(y-Xb)=wy-X^\top Xb.
\]

In the BED-derived workflow, these inputs are prepared in two
coordinated branches:

\begin{Shaded}
\begin{Highlighting}[]
\NormalTok{the same BED files + cls + rows + allele frequencies}
\NormalTok{    |}
\NormalTok{    +{-}{-}\textgreater{} make\_stats() {-}{-}{-}{-}{-}{-}\textgreater{} stats$wy, stats$ww, stats$yy}
\NormalTok{    |}
\NormalTok{    +{-}{-}\textgreater{} make\_sparseLD() {-}{-}{-}\textgreater{} Glist$sparseLD$prefix}
\NormalTok{                                      |}
\NormalTok{                                      v}
\NormalTok{                       stblr\_csr(stats = stats, Glist = Glist, ...)}
\end{Highlighting}
\end{Shaded}

\texttt{make\_stats()} supplies the phenotype cross-products and
diagonal marker cross-products. \texttt{make\_sparseLD()} supplies
retained off-diagonal marker correlations and records their prefix and
metadata in \texttt{Glist\$sparseLD}. During CSR model setup, the
sampler combines the stored correlations with \texttt{stats\$ww} to
recover the corresponding off-diagonal cross-product scale.

This split is important: the on-disk values are correlations, not raw
\(X_i^\top X_j\), and the diagonal cross-products are supplied
separately by \texttt{stats\$ww}.

\section{Matrix Represented by the Current
Writer}\label{matrix-represented-by-the-current-writer}

The current BED sparse-LD writer standardizes each selected marker using
its allele frequency, mean-imputes missing genotypes to zero on that
standardized scale, computes marker cross-products, and normalizes
off-diagonal entries to correlations:

\[
r_{ij}
=
\frac{x_i^\top x_j}
{\sqrt{(x_i^\top x_i)(x_j^\top x_j)}}.
\]

The writer retains an off-diagonal pair \(i<j\) when it passes the
configured distance rules and

\[
r_{ij}^2 \geq \texttt{r2_threshold}.
\]

The stored value is signed \(r_{ij}\), not \(r_{ij}^2\). The current
streamed artifact stores:

\begin{itemize}
\tightlist
\item
  only the upper triangle, \(j>i\);
\item
  no explicit diagonal entries;
\item
  an implicit diagonal correlation of one; and
\item
  only pairs passing the threshold and window rules.
\end{itemize}

Thus the disk artifact is a sparse, upper-triangle, correlation-like
representation of LD. It is not itself the complete \(X^\top X\) matrix.

When a CSR ST-BLR sampler opens the prefix, it reads the stored
correlation and constructs a symmetric in-memory neighbor structure. For
each retained pair it converts

\[
r_{ij}
\quad\longrightarrow\quad
\widetilde{X_i^\top X_j}
=
r_{ij}\sqrt{ww_iww_j}.
\]

The diagonal contribution remains \texttt{ww{[}i{]}}. This is why the
sparse-LD prefix and sufficient statistics must use compatible samples,
marker scaling, and marker order. The current shared-LD ST-BLR setup
also requires compatible \texttt{ww} across trait columns.

\section{What CSR Means}\label{what-csr-means}

Compressed sparse row format represents a sparse matrix without
allocating all \(m^2\) entries. Three arrays describe the retained
entries:

\begin{longtable}[]{@{}
  >{\raggedright\arraybackslash}p{(\linewidth - 2\tabcolsep) * \real{0.5000}}
  >{\raggedright\arraybackslash}p{(\linewidth - 2\tabcolsep) * \real{0.5000}}@{}}
\toprule\noalign{}
\begin{minipage}[b]{\linewidth}\raggedright
Array
\end{minipage} & \begin{minipage}[b]{\linewidth}\raggedright
Current conceptual role
\end{minipage} \\
\midrule\noalign{}
\endhead
\bottomrule\noalign{}
\endlastfoot
\texttt{row\_ptr} & Cumulative offsets delimiting each row's retained
entries; length \(m+1\). \\
\texttt{col\_idx} & Column index for every retained entry; length
\texttt{nnz}. \\
\texttt{values} & Value for every retained entry; length
\texttt{nnz}. \\
\end{longtable}

For zero-based row pointers, entries belonging to row \(i\) occupy
positions

\[
\texttt{row\_ptr}[i]
\quad\text{through}\quad
\texttt{row\_ptr}[i+1]-1.
\]

For example, consider the upper-triangle entries

\begin{Shaded}
\begin{Highlighting}[]
\NormalTok{row 0: (column 1, 0.40), (column 3, {-}0.20)}
\NormalTok{row 1: (column 3, 0.15)}
\NormalTok{row 2: no retained entries}
\NormalTok{row 3: no retained entries}
\end{Highlighting}
\end{Shaded}

Their CSR arrays are conceptually:

\begin{Shaded}
\begin{Highlighting}[]
\NormalTok{row\_ptr = [0, 2, 3, 3, 3]}
\NormalTok{col\_idx = [1, 3, 3]}
\NormalTok{values  = [0.40, {-}0.20, 0.15]}
\end{Highlighting}
\end{Shaded}

\texttt{row\_ptr} is always a cumulative pointer array. The
\texttt{one\_based} option of \texttt{sparseLD\_read\_CSR()} converts
returned \textbf{column indices} for convenient R inspection; it does
not convert the row pointers.

\section{Why Disk-Backed Sparse CSR Is
Used}\label{why-disk-backed-sparse-csr-is-used}

A dense \(m\times m\) LD matrix requires \(O(m^2)\) storage. This
becomes prohibitive for genome-scale marker sets even when most
scientifically useful LD relationships are local or sufficiently strong.

Sparse CSR reduces storage toward \(O(m+\mathrm{nnz})\), where
\texttt{nnz} is the number of retained off-diagonal upper-triangle
entries. Disk-backed storage provides two additional advantages:

\begin{itemize}
\tightlist
\item
  the resulting LD artifact can be reused by multiple CSR model fits;
  and
\item
  the complete sparse arrays do not have to be returned to or retained
  in R during construction.
\end{itemize}

\texttt{sparseLD\_stream\_CSR()} processes marker blocks, computes
candidate cross-products for block pairs, applies distance and \(r^2\)
rules, and writes retained entries incrementally. It still needs working
memory for decoded genotype blocks and block cross-products, but it
avoids constructing either a dense whole-genome LD matrix or the
complete sparse result in R.

\section{Stable User-Facing Contract}\label{stable-user-facing-contract}

The supported workflow treats the sparse-LD artifact as an opaque
package- managed object identified by a prefix:

\begin{Shaded}
\begin{Highlighting}[]
\NormalTok{ld\_prefix }\OtherTok{\textless{}{-}} \FunctionTok{file.path}\NormalTok{(output\_dir, }\StringTok{"chr1\_ld"}\NormalTok{)}

\FunctionTok{sparseLD\_stream\_CSR}\NormalTok{(}
  \AttributeTok{bed\_files =}\NormalTok{ Glist}\SpecialCharTok{$}\NormalTok{bedfiles[chr],}
  \AttributeTok{n =}\NormalTok{ Glist}\SpecialCharTok{$}\NormalTok{n,}
  \AttributeTok{cls =} \FunctionTok{list}\NormalTok{(cls),}
  \AttributeTok{out\_prefix =}\NormalTok{ ld\_prefix,}
  \AttributeTok{rows =} \ConstantTok{NULL}\NormalTok{,}
  \AttributeTok{af =} \FunctionTok{list}\NormalTok{(Glist}\SpecialCharTok{$}\NormalTok{af[[chr]][cls]),}
  \AttributeTok{r2\_threshold =} \FloatTok{0.001}\NormalTok{,}
  \AttributeTok{block\_size =} \DecValTok{1024}\NormalTok{,}
  \AttributeTok{nthreads =} \DecValTok{4}
\NormalTok{)}
\end{Highlighting}
\end{Shaded}

The stable usage expectations are:

\begin{itemize}
\tightlist
\item
  all files belonging to a prefix are created and retained together;
\item
  \texttt{sparseLD\_read\_CSR(ld\_prefix)} is used for inspection and
  validation;
\item
  \texttt{stblr\_csr(...,\ Glist\ =\ Glist)} and
  \texttt{stblr\_csr\_annot()} normally consume it through
  \texttt{Glist\$sparseLD\$prefix}; explicit \texttt{ld\_prefix} remains
  available for lower-level workflows;
\item
  marker order and construction settings are recorded externally; and
\item
  package functions handle file naming, binary types, index bases,
  diagonal convention, and conversion to sampler structures.
\end{itemize}

Analysis code should not assume that current suffixes, metadata keys, or
binary encodings will remain unchanged. Independent tools that manually
write these files depend on implementation details and can silently
become incompatible after package changes.

\section{\texorpdfstring{\texttt{ld\_prefix} and the Current Disk
Artifact}{ld\_prefix and the Current Disk Artifact}}\label{ld_prefix-and-the-current-disk-artifact}

\texttt{ld\_prefix} is a common path prefix, not a single file and not
an in-memory matrix. The current implementation derives several files
from it. For example, a prefix \texttt{path/chr1\_ld} currently
identifies:

\begin{Shaded}
\begin{Highlighting}[]
\NormalTok{path/chr1\_ld.row\_ptr.u64.bin}
\NormalTok{path/chr1\_ld.col\_idx.u32.0based.bin}
\NormalTok{path/chr1\_ld.values.f32.bin}
\NormalTok{path/chr1\_ld.meta.txt}
\end{Highlighting}
\end{Shaded}

These names expose current implementation details:

\begin{itemize}
\tightlist
\item
  row pointers are currently unsigned 64-bit integers;
\item
  column indices are currently unsigned 32-bit zero-based integers;
\item
  stored correlations are currently 32-bit floating-point values; and
\item
  metadata is currently a text file describing dimensions and settings.
\end{itemize}

The metadata currently records fields such as marker count,
\texttt{nnz}, triangle and diagonal conventions, value type,
normalization, threshold, distance limits, block size, thread count, and
used sample count.

Treat the complete prefix as one artifact. Do not mix files from
separate runs. Do not overwrite an important prefix when changing
samples, marker sets, allele frequencies, distance windows, or
thresholds. Prefer distinct, descriptive prefixes and retain
construction provenance beside them.

\section{Exported and Internal
Helpers}\label{exported-and-internal-helpers}

Current \texttt{NAMESPACE} and wrapper status distinguish the supported
public path from lower-level generated functions:

\begin{longtable}[]{@{}
  >{\raggedright\arraybackslash}p{(\linewidth - 4\tabcolsep) * \real{0.3333}}
  >{\raggedright\arraybackslash}p{(\linewidth - 4\tabcolsep) * \real{0.3333}}
  >{\raggedright\arraybackslash}p{(\linewidth - 4\tabcolsep) * \real{0.3333}}@{}}
\toprule\noalign{}
\begin{minipage}[b]{\linewidth}\raggedright
Helper
\end{minipage} & \begin{minipage}[b]{\linewidth}\raggedright
Current role
\end{minipage} & \begin{minipage}[b]{\linewidth}\raggedright
Current status
\end{minipage} \\
\midrule\noalign{}
\endhead
\bottomrule\noalign{}
\endlastfoot
\texttt{sparseLD\_stream\_CSR()} & Computes BED-derived sparse LD and
streams the disk-backed prefix. & Exported lower-level writer used by
\texttt{make\_sparseLD()}. \\
\texttt{sparseLD\_read\_CSR()} & Reads a written prefix into R for
inspection and validation. & Exported; user-facing reader. \\
\texttt{sparseLD\_to\_CSR()} & Computes the same upper-triangle sparse
representation and returns arrays in memory. & Generated lower-level
wrapper; not exported in the current \texttt{NAMESPACE}. \\
\texttt{sparseLD\_write\_CSR()} & Delegates to
\texttt{sparseLD\_stream\_CSR()} with an output prefix. & Generated
alias/internal wrapper; not exported in the current
\texttt{NAMESPACE}. \\
\end{longtable}

The lower-level wrappers can be useful during development, but they are
not the documented model handoff. In particular, \texttt{stblr\_csr()}
accepts \texttt{ld\_prefix}, not an R list returned by
\texttt{sparseLD\_to\_CSR()}.

The package also contains older or separate CSR conversion code for
other LD inputs. Its conventions should not be assumed to be identical
to the BED-streamed prefix described here.

\section{Sparse-LD Construction
Settings}\label{sparse-ld-construction-settings}

Sparse-LD settings affect memory, runtime, disk use, and the statistical
approximation. They should be selected deliberately and reported with
model results.

\subsection{Marker and Sample
Selection}\label{marker-and-sample-selection}

\begin{description}
\tightlist
\item[\texttt{bed\_files}]
Paths to one or more PLINK BED files. Their corresponding \texttt{cls}
entries define the concatenated model-marker order.
\item[\texttt{n}]
Number of individuals encoded by the BED source.
\item[\texttt{cls}]
A list of marker-column index vectors, one per BED file. The order
within and across these vectors becomes the sparse-LD row and column
order.
\item[\texttt{rows}]
Optional positive one-based sample indices. When supplied, LD is
computed from this selected sample. The same sample selection must
underlie the sufficient statistics used downstream.
\item[\texttt{af}]
Optional list of allele frequencies aligned to \texttt{cls}. When
omitted, the writer computes allele frequencies from the selected packed
genotypes. Supplied values must be finite and in \([0,1]\). Missing
genotypes are represented as zero after standardization, corresponding
to mean imputation on that scale.
\end{description}

\subsection{Threshold and Window
Settings}\label{threshold-and-window-settings}

\begin{description}
\tightlist
\item[\texttt{r2\_threshold}]
Minimum squared correlation for retaining an eligible pair. The writer
compares \(r_{ij}^2\) with this value but stores signed \(r_{ij}\). A
larger threshold produces a sparser approximation and may remove weak
but collectively relevant LD.
\item[\texttt{max\_distance\_variants}]
Maximum difference in model-order marker indices. Positive values
restrict retained pairs to this index window. A value of zero disables
this restriction in the current implementation.
\item[\texttt{max\_distance\_bp}]
Maximum base-pair separation when positions are supplied. Positive
values restrict retained pairs by physical distance. A value of zero
disables this restriction. The implementation evaluates forward position
differences, so \texttt{pos\_bp} should be aligned to marker order and
ordered consistently within the intended genomic segments.
\item[\texttt{pos\_bp}]
Optional list of marker base-pair positions aligned to \texttt{cls}.
When absent, the base-pair distance restriction cannot be applied, even
if \texttt{max\_distance\_bp} is positive.
\end{description}

A pair must pass every enabled window rule before its \(r^2\) is tested.
Disabling both distance restrictions permits all upper-triangle marker
pairs to be considered and can be computationally expensive. The
exported R wrapper fails closed in this case unless
\texttt{allow\_full\_ld\ =\ TRUE} is supplied explicitly.

\subsection{Computational Settings}\label{computational-settings}

\begin{description}
\tightlist
\item[\texttt{block\_size}]
Number of markers decoded and compared in a working block. Larger values
can improve matrix-multiplication efficiency but increase working memory
roughly through genotype-block and block-cross-product buffers. It
changes the construction strategy, not the intended retention rule.
\item[\texttt{nthreads}]
Requested OpenMP thread count. It affects construction performance and
may affect low-level floating-point execution details, but it is not an
LD threshold or model-prior parameter.
\end{description}

\subsection{Scaling}\label{scaling}

\texttt{sparseLD\_stream\_CSR()} does not expose a \texttt{scale}
argument. The current writer always constructs its correlation values
from allele-frequency-standardized markers. By contrast,
\texttt{bed\_xtx\_xty()} exposes \texttt{scale}. For the standard
BED-to-CSR handoff, use compatible standardized marker statistics:

\begin{Shaded}
\begin{Highlighting}[]
\NormalTok{stats }\OtherTok{\textless{}{-}} \FunctionTok{bed\_xtx\_xty}\NormalTok{(..., }\AttributeTok{af =}\NormalTok{ af, }\AttributeTok{scale =} \ConstantTok{TRUE}\NormalTok{)}
\FunctionTok{sparseLD\_stream\_CSR}\NormalTok{(..., }\AttributeTok{af =} \FunctionTok{list}\NormalTok{(af))}
\end{Highlighting}
\end{Shaded}

Using \texttt{bed\_xtx\_xty(...,\ scale\ =\ FALSE)} with the current
standardized correlation writer changes the scale relationship assumed
when the sampler combines \texttt{r\_ij} with \texttt{stats\$ww}. Such
combinations require explicit validation rather than being assumed
compatible.

\section{Streaming and Blocking
Behavior}\label{streaming-and-blocking-behavior}

The writer reads the selected BED markers into a packed representation,
then works through pairs of marker blocks. For each eligible block pair
it:

\begin{enumerate}
\def\labelenumi{\arabic{enumi}.}
\tightlist
\item
  decodes standardized marker values into floating-point working
  buffers;
\item
  computes raw block cross-products with matrix multiplication;
\item
  converts each eligible cross-product to correlation \(r_{ij}\);
\item
  applies distance and \(r^2\) retention rules; and
\item
  appends retained upper-triangle column indices and values to disk.
\end{enumerate}

After each row, it writes the cumulative retained-entry count to
\texttt{row\_ptr}. Consequently, the final row pointer equals
\texttt{nnz}.

Streaming avoids holding the final \texttt{col\_idx} and \texttt{values}
arrays in R or in one complete native in-memory result. It does not
eliminate all memory or runtime costs: the packed selected genotype
matrix, diagonal cross-products, decoded blocks, and block
cross-products still require resources. The number of candidate block
pairs can also become large when windows are wide or disabled.

\section{Index Bases and Diagonal
Convention}\label{index-bases-and-diagonal-convention}

Three conventions must be distinguished:

\begin{enumerate}
\def\labelenumi{\arabic{enumi}.}
\tightlist
\item
  \texttt{rows} and \texttt{cls} supplied from R use the documented
  R-facing indexing expected by their wrappers.
\item
  Current on-disk \texttt{col\_idx} entries are zero-based.
\item
  \texttt{sparseLD\_read\_CSR(prefix,\ one\_based\ =\ TRUE)} returns
  one-based column indices for convenient R inspection.
\end{enumerate}

The reader's \texttt{one\_based} argument changes only the returned
\texttt{col\_idx} vector. \texttt{row\_ptr} remains a zero-based
cumulative pointer array whose first value is zero and whose final value
is \texttt{nnz}.

The current streamed artifact omits diagonal entries and reports
\texttt{diag\ =\ "implicit\_1"}. Downstream ST-BLR samplers use the
supplied \texttt{stats\$ww} for diagonal \(X_i^\top X_i\) contributions.
Users inspecting the sparse arrays should not mistake absence of
explicit diagonal entries for a zero diagonal.

Package functions validate and handle these conventions. Manual index
conversion or insertion of diagonal values risks corrupting the expected
contract.

\section{Marker-Order, Sample, and Scaling
Invariants}\label{marker-order-sample-and-scaling-invariants}

Sparse arrays contain positional indices, not biological marker
identifiers. The most important cross-layer invariant is:

\begin{Shaded}
\begin{Highlighting}[]
\NormalTok{Glist$rsids[[chr]][cls]}
\NormalTok{    == Glist$rsidsLD[[chr]]}
\NormalTok{    == rows of stats$wy and stats$ww}
\NormalTok{    == rows and columns of sparse LD}
\NormalTok{    == rows of an annotation matrix A}
\NormalTok{    == rows of fit$bm and fit$dm}
\end{Highlighting}
\end{Shaded}

Establish and retain marker order explicitly:

\begin{Shaded}
\begin{Highlighting}[]
\NormalTok{marker\_id }\OtherTok{\textless{}{-}}\NormalTok{ Glist}\SpecialCharTok{$}\NormalTok{rsidsLD[[chr]]}
\NormalTok{cls }\OtherTok{\textless{}{-}} \FunctionTok{match}\NormalTok{(marker\_id, Glist}\SpecialCharTok{$}\NormalTok{rsids[[chr]])}

\FunctionTok{stopifnot}\NormalTok{(}
  \SpecialCharTok{!}\FunctionTok{anyNA}\NormalTok{(cls),}
  \FunctionTok{identical}\NormalTok{(Glist}\SpecialCharTok{$}\NormalTok{rsids[[chr]][cls], marker\_id)}
\NormalTok{)}
\end{Highlighting}
\end{Shaded}

For multiple BED files, the concatenation of \texttt{cls{[}{[}1{]}{]}},
\texttt{cls{[}{[}2{]}{]}}, and so on defines the CSR order.
Allele-frequency and position lists must follow the same per-file and
concatenated order.

The same selected sample must underlie LD and sufficient statistics. If
\texttt{rows} is used, apply it consistently. If allele frequencies are
supplied, they must correspond to that construction and marker order. If
annotations are later used, align their rows before fitting:

\begin{Shaded}
\begin{Highlighting}[]
\FunctionTok{rownames}\NormalTok{(A) }\OtherTok{\textless{}{-}}\NormalTok{ marker\_id}
\FunctionTok{stopifnot}\NormalTok{(}\FunctionTok{identical}\NormalTok{(}\FunctionTok{rownames}\NormalTok{(A), marker\_id))}
\end{Highlighting}
\end{Shaded}

A permutation error can produce dimensionally valid CSR files and model
fits while assigning LD or annotations to the wrong markers. Dimension
checks alone cannot detect this error.

\section{Reading and Validating a
Prefix}\label{reading-and-validating-a-prefix}

Use the exported reader to inspect a prefix without manually decoding
the binary files:

\begin{Shaded}
\begin{Highlighting}[]
\NormalTok{ld }\OtherTok{\textless{}{-}} \FunctionTok{sparseLD\_read\_CSR}\NormalTok{(ld\_prefix, }\AttributeTok{one\_based =} \ConstantTok{TRUE}\NormalTok{)}

\FunctionTok{stopifnot}\NormalTok{(}
\NormalTok{  ld}\SpecialCharTok{$}\NormalTok{nrow }\SpecialCharTok{==} \FunctionTok{length}\NormalTok{(cls),}
\NormalTok{  ld}\SpecialCharTok{$}\NormalTok{ncol }\SpecialCharTok{==} \FunctionTok{length}\NormalTok{(cls),}
  \FunctionTok{length}\NormalTok{(ld}\SpecialCharTok{$}\NormalTok{row\_ptr) }\SpecialCharTok{==} \FunctionTok{length}\NormalTok{(cls) }\SpecialCharTok{+} \DecValTok{1}\NormalTok{L,}
  \FunctionTok{length}\NormalTok{(ld}\SpecialCharTok{$}\NormalTok{col\_idx) }\SpecialCharTok{==} \FunctionTok{length}\NormalTok{(ld}\SpecialCharTok{$}\NormalTok{values),}
  \FunctionTok{tail}\NormalTok{(ld}\SpecialCharTok{$}\NormalTok{row\_ptr, }\DecValTok{1}\NormalTok{) }\SpecialCharTok{==} \FunctionTok{length}\NormalTok{(ld}\SpecialCharTok{$}\NormalTok{values),}
\NormalTok{  ld}\SpecialCharTok{$}\NormalTok{index\_base }\SpecialCharTok{==} \DecValTok{1}\NormalTok{,}
  \FunctionTok{identical}\NormalTok{(ld}\SpecialCharTok{$}\NormalTok{diag, }\StringTok{"implicit\_1"}\NormalTok{)}
\NormalTok{)}
\end{Highlighting}
\end{Shaded}

The current reader checks that required metadata dimensions exist, reads
the expected byte counts, verifies that row pointers begin at zero and
end at \texttt{nnz}, requires nondecreasing row pointers, checks
column-index bounds, and rejects non-finite values.

Additional analysis-level validation should confirm:

\begin{enumerate}
\def\labelenumi{\arabic{enumi}.}
\tightlist
\item
  \texttt{cls} contains no missing marker matches.
\item
  LD dimensions equal the intended marker count.
\item
  \texttt{row\_ptr} has length \(m+1\).
\item
  \texttt{col\_idx} and \texttt{values} have equal length.
\item
  the final row pointer equals \texttt{nnz}.
\item
  returned \texttt{index\_base} matches how \texttt{col\_idx} is being
  inspected.
\item
  the upper-triangle and implicit-diagonal conventions are understood.
\item
  values are finite signed correlations and remain within plausible
  correlation bounds.
\item
  \texttt{n\_used}, thresholds, windows, and marker count match the
  intended run.
\item
  \texttt{stats\$wy}, every \texttt{stats\$ww} vector, annotations, and
  posterior rows use exactly the same marker order.
\end{enumerate}

The reader can validate structural properties but cannot prove that the
BED source, allele frequencies, sample rows, phenotypes, annotations, or
marker identities are scientifically aligned. Retain that provenance
separately.

\section{\texorpdfstring{How \texttt{stblr\_csr()} Consumes the
Prefix}{How stblr\_csr() Consumes the Prefix}}\label{how-stblr_csr-consumes-the-prefix}

The public model handoff is short:

\begin{Shaded}
\begin{Highlighting}[]
\NormalTok{fit }\OtherTok{\textless{}{-}} \FunctionTok{stblr\_csr}\NormalTok{(}
  \AttributeTok{stats =}\NormalTok{ stats,}
  \AttributeTok{Glist =}\NormalTok{ Glist,}
  \AttributeTok{scheduled =} \ConstantTok{FALSE}
\NormalTok{)}
\end{Highlighting}
\end{Shaded}

During setup, the current ST-BLR CSR implementation:

\begin{enumerate}
\def\labelenumi{\arabic{enumi}.}
\tightlist
\item
  checks that metadata marker count matches the model marker count;
\item
  reads the zero-based upper-triangle arrays;
\item
  validates row pointers, column bounds, finite correlations, and
  positive \texttt{stats\$ww};
\item
  ignores any explicit self-pairs if encountered;
\item
  converts retained \(r_{ij}\) to \(r_{ij}\sqrt{ww_iww_j}\); and
\item
  creates a symmetric in-memory neighbor structure for residual updates.
\end{enumerate}

For marker effect \(b_i\), the diagonal residual contribution uses
\(ww_i b_i\), while retained off-diagonal neighbors use the converted
cross-product values. Missing off-diagonal entries are treated as zero
in the sparse approximation.

Annotation-aware CSR wrappers use the same LD prefix and marker-order
contract. Their differences concern prior structure, not CSR
construction.

\section{Scheduled Sparse Updates}\label{scheduled-sparse-updates}

\texttt{stblr\_csr(...,\ scheduled\ =\ TRUE)} uses a scheduled
marker-update implementation that can revisit likely null markers less
frequently and wake retained LD neighbors after an effect change.
Settings such as \texttt{full\_sweep\_every},
\texttt{candidate\_threshold}, \texttt{candidate\_lifetime},
\texttt{null\_skip\_base}, \texttt{null\_skip\_max}, and LD-neighbor
wake-up controls affect the sampler's computational schedule.

These controls are separate from CSR construction settings:

\begin{itemize}
\tightlist
\item
  \texttt{r2\_threshold} and distance windows determine which LD edges
  exist;
\item
  scheduling controls determine when markers and retained neighbors are
  reconsidered during MCMC.
\end{itemize}

Both can affect finite-run behavior and performance. The statistical
prior is unchanged by choosing a scheduled path, but aggressive sparse
scheduling must be checked for adequate mixing and sensitivity. See
\href{single_trait_models.html}{Single-trait models} for the scheduling
controls.

\section{Sparse LD as an
Approximation}\label{sparse-ld-as-an-approximation}

The sparse matrix replaces unretained off-diagonal relationships with
zero. Approximation enters through:

\begin{itemize}
\tightlist
\item
  \(r^2\) thresholding;
\item
  marker-index and physical-distance windows;
\item
  finite-sample LD estimation;
\item
  allele-frequency estimation and missing-genotype handling;
\item
  32-bit storage of retained correlation values; and
\item
  any difference between the sample or scaling used for LD and
  sufficient statistics.
\end{itemize}

Increasing \texttt{r2\_threshold} or narrowing windows usually reduces
\texttt{nnz}, runtime, and disk use, but can remove LD needed to
represent the conditional marker likelihood adequately. Lower thresholds
and wider windows preserve more LD at greater computational cost. There
is no universally correct threshold; its adequacy depends on marker
density, LD structure, sample size, phenotype architecture, and
inferential target.

\section{BED-Versus-CSR Comparison and
Reproducibility}\label{bed-versus-csr-comparison-and-reproducibility}

When both paths use the same BED source, markers, samples, phenotype,
standardization, and priors, \texttt{stblr\_bed()} and
\texttt{stblr\_csr()} should usually show qualitatively comparable
posterior patterns. Exact equality is not expected because:

\begin{itemize}
\tightlist
\item
  the CSR path omits unretained LD;
\item
  BED and CSR samplers maintain residual information differently;
\item
  update scheduling and initialization can differ;
\item
  retained correlations use finite precision; and
\item
  MCMC introduces Monte Carlo variation.
\end{itemize}

A comparison should examine PIP rankings, posterior mean effects,
aggregate variance summaries, and sensitivity to thresholds rather than
equality of individual samples.

For reproducibility, retain:

\begin{itemize}
\tightlist
\item
  immutable identifiers for BED inputs and selected samples;
\item
  marker IDs in exact CSR order;
\item
  \texttt{cls}, \texttt{rows}, allele frequencies, and positions;
\item
  sufficient-statistic dimensions and scaling choices;
\item
  all files belonging to \texttt{ld\_prefix};
\item
  the writer's returned summary and metadata;
\item
  package version and platform information; and
\item
  model priors, MCMC settings, scheduling controls, and seeds.
\end{itemize}

The current workflow scripts use demonstration-scale settings.
Production analyses require sensitivity checks for LD construction and
sampler settings, longer chains, and convergence diagnostics.

\section{Related Notes and Sources}\label{related-notes-and-sources}

\begin{itemize}
\tightlist
\item
  \href{sparse_ld_and_bed_workflows.html}{Sparse-LD and BED workflows}
\item
  \href{data_representations.html}{Data representations}
\item
  \href{single_trait_models.html}{Single-trait models}
\item
  \href{workflows_and_outputs.html}{Workflows and outputs}
\item
  \href{technical_summary_statistics.html}{Technical summary statistics}
\item
  \href{https://github.com/psoerensen/sblr/blob/master/examples/workflows/sparse_ld_bed_workflow.R}{Complete
  sparse-LD/BED workflow on GitHub}
\item
  \href{https://github.com/psoerensen/sblr/blob/master/examples/workflows/annotation_based_models.R}{Complete
  annotation workflow on GitHub}
\end{itemize}




\end{document}
